\chapter{Introdução}
\label{cap:introducao}

\notaguia{Primeira redação, a partir da ``tese em uma página'' do guia de
leitura, da abertura da visualização e do escopo original. Números aqui são
de enquadramento (período, unidades, atos); os resultados quantitativos
moram no Capítulo~\ref{cap:resultados}.}

\section{Contexto e problema}
\label{sec:intro-contexto}

O Cerrado é a fronteira ativa da agricultura brasileira, e Goiás está no
coração do bioma. A série anual do MapBiomas (Coleção 10.1) permite observar,
ano a ano e pixel a pixel, o que cada fração do território goiano foi entre
1985 e 2024 --- pastagem, lavoura, vegetação natural, área urbana --- e,
portanto, reconstituir quatro décadas de transformação da paisagem com uma
resolução que era impensável há pouco tempo.

O que essa série mostra para Goiás não é uma história simples de avanço de
fronteira. O estado já entra na série amplamente ocupado: a pastagem que
domina a paisagem de 1985 é herança das frentes de ocupação anteriores. O
que as quatro décadas seguintes registram é uma \emph{reorganização}
profunda e espacialmente estruturada desse território já ocupado: no Sul, a
pastagem cede lugar à lavoura e a produção se intensifica; no Norte, a
fronteira segue ativa, convertendo vegetação natural em pasto. O resultado
agregado é que lavoura, pastagem e rebanho deslocam-se gradualmente para o
norte do estado --- o fenômeno que este trabalho chama de ``marcha ao
norte''.

O problema de pesquisa é duplo: descrever esse padrão com precisão --- ele
existe, é robusto, em que ritmo avança? --- e explicar o que o coordena.
Três famílias de explicação concorrem: uma região empurrando a outra (o
deslocamento indireto de uso da terra); forças de mercado comuns atuando
sobre um gradiente natural de aptidão; e o próprio encolhimento do que resta
exposto à conversão. Distinguir entre elas exige mais que
mapas: exige o cruzamento sistemático da série de uso da terra com as
variáveis socioeconômicas --- produto agropecuário, crédito rural, rebanho,
preços de commodities, câmbio --- e um desenho analítico capaz de rejeitar
as explicações que não se sustentam.

\section{Pergunta e objeto}
\label{sec:intro-pergunta}

A pergunta central é: \emph{como e por que o uso da terra mudou em Goiás
entre 1985 e 2024?} --- e, em sua forma mais aguda: \emph{o que coordena a
reorganização espacial da produção agropecuária no estado?}

O objeto empírico são os 246 municípios goianos, observados anualmente ao
longo de 40 anos. Para as análises longitudinais, os municípios são
agregados em 166 Áreas Mínimas Comparáveis, malha que mantém o território
constante diante das emancipações municipais do período
(Capítulo~\ref{cap:metodologia}). Sobre essa base, a série de uso e
cobertura da terra é cruzada com o painel socioeconômico municipal.

A investigação partiu de uma hipótese clássica --- a pastagem como estágio
intermediário da ocupação (vegetação $\to$ pasto $\to$ lavoura), com a
lavoura do Sul empurrando a pecuária para o Norte do próprio estado --- e a
trajetória da pesquisa foi, em boa medida, o teste dessa hipótese. Parte
dela se confirmou; a parte mais tentadora, não. O deslocamento causal de uma
região sobre a outra, testado formalmente, não se sustentou, e o que emergiu
em seu lugar foi a leitura que estrutura este texto: duas dinâmicas
paralelas --- intensificação no Sul, fronteira no Norte --- movidas pelo
mesmo motor econômico sobre um gradiente de aptidão, e limitadas por um teto
de oferta de terra convertível. Vários dos resultados centrais do trabalho
são, assim, resultados nulos documentados: afirmações sobre o que os dados
mostram que \emph{não} acontece.

\section{Justificativa}
\label{sec:intro-justificativa}

A lacuna científica é de escala e de desenho. A literatura empírica da
fronteira do Cerrado opera majoritariamente na escala do bioma e das suas
grandes regiões --- com o MATOPIBA, acrônimo que designa a fronteira agrícola
formada por Maranhão, Tocantins, Piauí e Bahia, como caso paradigmático --- e a
literatura de mudança indireta de uso da terra (iLUC, do inglês \emph{indirect
land use change}) estabelece o fenômeno entre regiões e
biomas \cite{Lapola2010, Arima2011}. O canal \emph{intra-estadual} --- uma
região de um mesmo estado empurrando a fronteira da outra --- raramente é
isolado e testado, e as séries municipais longas costumam ignorar o problema
das fronteiras administrativas móveis. Este trabalho contribui exatamente
nesse ponto: alta resolução espacial, território constante e um teste formal
do canal de deslocamento local (Capítulo~\ref{cap:referencial}).

A relevância ambiental é direta: o Cerrado remanescente de Goiás
concentra-se no Norte do estado --- precisamente a região para onde a
conversão migrou. Compreender o que move essa conversão, e em que ritmo,
é condição para dimensionar a exposição da vegetação que resta. Para
políticas públicas, a distinção importa na prática: se a conversão no Norte
fosse efeito-dominó da lavoura do Sul, conter a lavoura conteria a
fronteira; se, como os resultados sugerem sem ainda estabelecer, ambas
respondem a um drive de mercado comum, as alavancas eficazes são outras ---
e agem sobre o estado inteiro.

Por fim, no espírito do mestrado profissional, o trabalho gera produtos de
uso direto: um conjunto reproduzível de pipelines de coleta e análise sobre
dados públicos e uma visualização interativa que torna os achados
navegáveis por não especialistas.

\section{Objetivos}
\label{sec:intro-objetivos}

O objetivo geral é caracterizar e explicar a reorganização espacial do uso e
cobertura da terra em Goiás entre 1985 e 2024, em sua relação com fatores
socioeconômicos.

Os objetivos específicos espelham as quatro frentes da investigação:

\begin{enumerate}
  \item caracterizar o padrão espaço-temporal da conversão de terras e
        verificar a existência, o ritmo e a robustez do deslocamento da
        fronteira para o norte;
  \item identificar os mecanismos locais que sustentam o padrão ---
        intensificação e substituição pasto$\to$lavoura no Sul,
        extensificação sobre vegetação natural no Norte;
  \item testar formalmente a hipótese de deslocamento intra-estadual e
        avaliar o papel de drivers macroeconômicos comuns, com destaque
        para o câmbio;
  \item examinar a desaceleração recente da conversão e a hipótese do teto
        de oferta de terra convertível no Sul do estado.
\end{enumerate}

\section{Organização do documento}
\label{sec:intro-organizacao}

O texto está organizado em seis capítulos, precedidos por um memorial. O
Capítulo~\ref{cap:referencial} situa a investigação na literatura da
fronteira agrícola, do deslocamento indireto de uso da terra, da
intensificação e dos drivers macroeconômicos. O
Capítulo~\ref{cap:metodologia} apresenta as fontes de dados, as unidades de
análise, a construção do painel e as decisões metodológicas críticas ---
inclusive o papel da inteligência artificial como instrumento de pesquisa.
O Capítulo~\ref{cap:resultados} apresenta os resultados preliminares em duas
camadas: a periodização dos 40 anos em três atos (1985--2000, 2001--2019,
2020--2024) e as quatro frentes de evidência da marcha ao norte. O
Capítulo~\ref{cap:discussao} posiciona os achados na literatura e consolida
alcance e limites, e o Capítulo~\ref{cap:cronograma} apresenta o cronograma
até a defesa.
